As synthetic text pervades web corpora, detect-and-reweight methods that downweight likely-synthetic samples during training have emerged to prevent model collapse, yet when they actually help remains unclear.
We investigate this question through controlled continued-pretraining experiments (Qwen3-0.6B, 132M tokens, 5 seeds) and identify three effectiveness conditions.
First, detect-and-remove (D\&R) outperforms all reweighting methods in our setting; epoch-matched controls attribute over 95\% of the advantage to multi-epoch exposure on filtered data, identifying the mechanism behind D\&R's effectiveness.
Second, the detector must provide both sufficient discrimination \emph{and} good calibration; high-AUC detectors with poor calibration degrade performance below uniform training.
Third, minimum intervention, applying the smallest reweighting exponent that meets a quality floor, outperforms unconstrained weight maximization.
As an illustrative application of the minimum-intervention principle, we derive a dual-constrained adaptive exponent $\bstar$ that improves over uniform training at contamination ratios $\rho \geq 0.4$, reaching conventional significance after multiple-comparison correction, while requiring no per-model tuning.
These findings apply within our self-contamination setup at 0.6B scale; generalization to cross-model contamination and larger models remains to be established\footnote{Code and data are available at \url{https://anonymous.4open.science/r/provenance_weight_training-F72D/}.}.
